\section{Sequencer}
\label{sec:sequencer}
The Sequencer is the entry point for L2 transactions. It is responsible for ingesting, validating, and grouping transactions into batches before forwarding them to the execution stage. Strict architectural rules dictate that the sequencer interacts with its database via a defined \texttt{Storage} trait and keeps business logic separate from the \texttt{src/bin/} entry point.

\begin{itemize}
    \item \textbf{Language:} Rust.
    \item \textbf{API Interface:} Exposes a JSON-RPC Server to receive transactions from users and the Python Workload Generator.
    \item \textbf{Transaction Flow:} Incoming transactions pass through a Validity Checker that interfaces with a pessimistic state cache (verifying nonces and balances). Valid transactions are pushed to a Normal Tx Pool. 
    \item \textbf{Batch Formation:} A background Batch Orchestrator pulls from these pools to form batches based on configured triggers. The primary triggers are reaching a maximum batch size (\texttt{max\_batch\_size}) or exceeding a defined time interval (\texttt{timeout\_interval\_ms}).
    \item \textbf{Scheduling Policies:} The sequencer uses the Strategy pattern to support different ordering policies, specifically First-Come-First-Served (FCFS), Fee-Priority, Time-Boost, Fair BFT, and BlobPacking.
    \item \textbf{Output:} Completed batches are dispatched to the Executor via a gRPC stream. Configuration is managed via TOML files (e.g., \texttt{sequencer.template.toml}), which are processed via \texttt{envsubst} during container startup (notably requiring explicit bash environment variables since default fallbacks are unsupported).
\end{itemize}
